\chapter{Agentes especializados y meta-agente}
\label{cap:agentes}

Este capítulo describe el corazón de aprendizaje del sistema: tres agentes especializados, cada uno
centrado en una familia de factores, y un meta-agente que los combina. Se explica qué predice cada
agente, con qué modelo, cómo se entrena para no incurrir en \emph{lookahead} y por qué se pasó de un
modelo lineal a árboles con \emph{gradient boosting}. Las cifras de \rankic{} que aparecen en este
capítulo se citan como \textbf{evidencia que justificó una decisión de método} ---por qué se cambió
de modelo, por qué se pondera de una forma y no de otra---, no como resultado final del sistema; los
resultados del estudio se recogen en el capítulo~\ref{cap:resultados}.

\section{Agentes especializados por familia de factores}

El sistema no entrena un único modelo sobre todas las variables, sino \textbf{tres agentes}
especializados, cada uno alineado con una prima de factor clásica de las revisadas en el
capítulo~\ref{cap:estado_arte}:

\begin{itemize}
  \item \textbf{Calidad}: rentabilidad, estabilidad y solidez del negocio.
  \item \textbf{Momentum}: fuerza relativa reciente del precio.
  \item \textbf{Valor}: baratura de la acción respecto a sus fundamentales.
\end{itemize}

La especialización tiene una doble justificación. Económicamente, cada agente encarna una hipótesis
de inversión con respaldo en la literatura, lo que hace el sistema \emph{explicable}: se puede decir
por qué una acción está bien valorada según cada estilo. Metodológicamente, separar las familias
permite medir el aprendizaje de cada una por separado y combinarlas después de forma controlada, en
lugar de dejar que un único modelo mezcle señales de naturaleza distinta de manera opaca. Todos los
factores se calculan como rangos dentro del corte transversal de cada fecha, de modo que lo que el
sistema aprende y evalúa es un \emph{orden} de activos, coherente con la métrica \rankic{}.

\section{Del modelo lineal a LightGBM}

El sistema partió de agentes \textbf{lineales} (una regresión sobre los factores) que predecían el
exceso de retorno. Esa base se midió con rigor ---entrenamiento \emph{walk-forward} \emph{point-in-time},
\rankic{} fuera de muestra--- y se probaron sobre ella múltiples mejoras: neutralización por sector,
distintos tratamientos de la etiqueta, \emph{momentum} y descomposición de fundamentales, y variables
de régimen de mercado. El resultado fue consistente: el \rankic{} fuera de muestra del enfoque lineal
se mantuvo \textbf{cercano a cero} en todas las variantes, con el mejor apenas en el entorno de
\(+0.001\) a \(+0.002\).

El diagnóstico de por qué se desarrolla en \textcite{gu2020empirical} y se confirmó aquí: un modelo
lineal \textbf{promedia a cero} las interacciones entre factores, de modo que añadir variables no
crea señal donde no la hay ---solo añade lugares donde sobreajustar---. Se concluyó que el techo del
enfoque lineal estaba en \rankic{} \(\approx 0\), y se adoptó \textbf{LightGBM} \parencite{ke2017lightgbm},
un modelo de árboles con \emph{gradient boosting} capaz de representar las interacciones no lineales
que el lineal no captura, con el objetivo \texttt{rank\_regression} (una regresión sobre el percentil
transversal del retorno futuro, alineada con el \rankic{}).

El cambio produjo una mejora \emph{relativa} clara ---en las comparaciones directas el \rankic{} del
meta con LightGBM casi duplicó al del lineal--- y robusta frente a la semilla aleatoria. Es
importante subrayar el matiz honesto que acompaña a este dato y que se retoma en los resultados: una
mejora relativa sobre una base próxima a cero sigue siendo, en términos absolutos, una señal muy
débil frente al listón que la industria considera útil (del orden de \(0.03\)--\(0.05\) de \rankic{}
sostenido). La mejora justifica \emph{el modelo}; no equivale por sí sola a haber aprendido de forma
significativa.

\section{Entrenamiento \emph{walk-forward}}

Cada agente se entrena de forma \textbf{\emph{walk-forward}}: en cada reentreno solo se usa historia
anterior cuya etiqueta ya se ha \emph{realizado}. Con la notación de la sección~\ref{sec:notacion},
una observación de fecha \(\tsnap\) tiene su etiqueta en \(\tsnap + \hlabel\); esa observación solo
entra en un entrenamiento cuyo corte sea posterior a \(\tsnap + \hlabel\). Así se impide que el
modelo se entrene sobre un resultado que en la práctica todavía no se conocería, cerrando la vía de
fuga temporal por el lado de la etiqueta ---complementaria a la observabilidad de los datos de
entrada del capítulo~\ref{cap:metodologia}---.

\section{Meta-agente: ponderar por aprendizaje, no por rentabilidad}

Los tres agentes se combinan en una única señal ---el \emph{meta-score}--- mediante un meta-agente.
La combinación no es una media simple: por defecto pondera cada agente por su \textbf{\rankic{}
reciente}, de modo que pesa más el agente que mejor ha estado ordenando los activos fuera de muestra
en el pasado próximo. En las comparaciones del proyecto, esta ponderación por \rankic{} superó a la
equiponderación, lo que indica que la combinación aporta y que la calidad reciente de cada agente es
informativa de su calidad inmediata.

Del \emph{meta-score} se deriva su percentil transversal (el \emph{meta-rank}), y es ese rango lo que
consume la cartera (capítulo~\ref{cap:diseno_experimental}). Coherentemente con la métrica principal
del trabajo, el \rankic{} se mide sobre el \emph{meta-score} final ---la señal que realmente opera---
y no sobre el promedio de los agentes individuales, y se reporta \textbf{separado} de cualquier medida
de rentabilidad: el aprendizaje se juzga por cómo ordena, y la rentabilidad es una consecuencia
posterior que se analiza aparte.
